\documentclass[11pt]{article}

%%%%%%%%%%%%%%Include Packages%%%%%%%%%%%%%%%%%%%%%%%%%%
\usepackage{xcolor}
\usepackage{mathtools}
\usepackage[a4paper, total={6in, 8in}, margin=1.25in]{geometry}
\usepackage{amsmath}
\usepackage{amssymb}
\usepackage{paralist}
\usepackage{rsfso}
\usepackage{amsthm}
\usepackage{wasysym}
\usepackage[inline]{enumitem}   
\usepackage{hyperref}
\usepackage{tocloft}
\usepackage{titlesec}
\usepackage{colortbl}
\usepackage{wrapfig}
\usepackage{pdfpages}
%%%%%%%%%%%%%%%%%%%%%%%%%%%%%%%%%%%%%%%%%%%%%%%%%%%%%%%%



\begin{document}
\Large \noindent \textbf{Physics 5A Discussion (Week 4)}\\
\normalsize
Department of Physics and Astronomy, UCLA\\
\noindent\rule{8cm}{0.4pt}\\

Recall the rocket we discussed in week 1 and week 2, but you don't need anything you calculated in those worksheets as these are completely separate problems. \textbf{We set up the coordinate system such that the rocket starts its motion at the origin, $x$-axis being the horizontal motion, $y$-axis being the vertical motion}. Here is some space for you to doodle if you want :)\\ 
\hfill\newline\hfill\newline\hfill\newline
\hfill\newline\hfill\newline\hfill\newline

\noindent(a) From time $t = 0$ to $t=0.2$ seconds the rocket moves to the left at a constant velocity. Then the boost (accelerating it to the right) is turned on, which slows down the leftward motion, until the rocket stops at $t=0.37$ seconds, at a distance of $0.285$ meters to the left of the origin. Write down the displacement vector of the rocket from $t=0$ to $t=0.37$ seconds. Write it in component form, that is $\Delta\vec{x} = (x_{\text{f}}-x_{\text{i}}\, ,\, y_{\text{f}}-y_{\text{i}})$.\\
\hfill\newline\hfill\newline\hfill\newline
\hfill\newline\hfill\newline\hfill\newline

\noindent(b) At time $t=0.37$ seconds, the rocket starts moving to the right. At time $t=0.54$ seconds, the rocket ends at a distance $0.2$ meters to the left of the origin. Write down the displacement vector of the rocket from $t=0.37$ to $t=0.54$ seconds. (Hint: Do ``final minus initial".)\\
\hfill\newline\hfill\newline\hfill\newline
\hfill\newline\hfill\newline\hfill\newline

\noindent(c) Write down the displacement vector for its journey from $t=0$ to $t=0.54$ seconds. This is different from part (b). Remember, "final minus initial".\\

\newpage
\noindent(d) Now we consider a two-step process:
\begin{enumerate}
\item From $t=0.54$ to $t=0.80$ seconds, the rocket goes from $0.2$ meters to the left of the origin back to the origin. 
\item From $t = 0.80$ to $t = 16.40$ seconds, the rocket goes in the "up-hill" direction (I called it the $x'$-direction in the previous worksheet), making an angle $30^\circ$ above the horizon, and it travels $200$ meters in that direction. 
\end{enumerate}
Find the component form of the displacement vector in each of the two processes. That is, the displacement from $t=0.54$ to $0.80$ seconds, and the displacement from $t=0.80$ to $16.40$ seconds.\\
\hfill\newline\hfill\newline\hfill\newline
\hfill\newline\hfill

\noindent(e) Find the total displacement from $t=0.54$ to $16.4$ seconds. (Hint: You will need to add vectors.)\\
\hfill\newline\hfill\newline\hfill\newline
\hfill\newline\hfill

\noindent(f) Find the average velocity from $t = 0.54$ to $16.4$ seconds, in component form.\\
\hfill\newline\hfill\newline\hfill\newline
\hfill\newline\hfill

\noindent(g) Find the total displacement from $t=0$ to $16.4$ seconds, first write it in component form, then compute its magnitude and angle. \\
\hfill\newline\hfill\newline\hfill\newline
\hfill\newline\hfill

\noindent(h) Find the total distance traveled by the rocket from $t = 0$ to $ 16.4$ seconds.\\ 
\hfill\newline\hfill\newline\hfill\newline
\hfill\newline\hfill

\noindent(i) Find the average speed of the rocket from $t = 0$ to $ 16.4$ seconds. 

\newpage
\textbf{Solution to (a).} The rocket ends at a point directly to the left of the origin, so we expect to see only $x$-component in its displacement. 
\begin{align*}
\Delta \vec{x} = (x_{\text{f}}-x_{\text{i}}\, ,\, y_{\text{f}}-y_{\text{i}}) = (-0.285 - 0,\ 0-0) = (-0.285,\ 0)
\end{align*}
with unit of meters.\\

\textbf{Solution to (b).} Starting position at $t=0.37$ seconds is $(-0.285,\ 0)$, and ending position is $(-0.2,\ 0)$ at $t = 0.54$, so
\begin{align*}
\Delta \vec{x} = (-0.2,\ 0) - (-0.285,\ 0) = (0.085, \ 0)\,,
\end{align*}
with unit of meters. The $x$-component is positive because it moves to the right from $t=0.37$ to $t=0.54$ seconds.\\

\textbf{Solution to (c).}  Starting position at $t=0$ seconds is at the origin $(0,\ 0)$, and ending position is $(-0.2,\ 0)$ at $t = 0.54$ seconds, so
\begin{align*}
\Delta \vec{x} = (-0.2,\ 0) - (0,\ 0) = (-0.2, \ 0)\,,
\end{align*}
with unit of meters. The $x$-component is negative because it ends at a point to the left of where it started.\\

\textbf{Solution to (d).} First, for displacement from $t=0.54$ to $0.80$ seconds, it starts at $(-0.2,\ 0)$, and ends at $(0,\ 0)$, so 
\begin{align*}
\Delta \vec{x}_1 = (0,\ 0) - (-0.2,\ 0) = (0.2, \ 0)\,
\end{align*}
with unit of meters. Then, for displacement from $t=0.80$ to $16.40$ seconds, it starts at $(0,\ 0)$, and ends at a point $\vec{p}$ that we need to calculate. The length that the rocket has traveled is $200$ meters, this is the hypotenuse, and the angle is $30^\circ$ as mentioned in the problem, so the vector $\vec{p}$ has $x$- and $y$-component that we can calculate using trigonometry, 
\begin{align*}
p_x = 200 \cos(30^\circ) = 173.2\,\text{m}\,,\qquad
p_y = 200 \sin(30^\circ) = 100 \,\text{m}\,.
\end{align*}
So the displacement for the second step is
\begin{align*}
\Delta\vec{x}_2 = ( 173.2,\ 100) -(0,0) =( 173.2,\ 100)
\end{align*}
with unit of meters. \\

\textbf{Solution to (e).} The displacement from $t=0.8$ to $16.4$ seconds can be calculated in two ways. First way by adding,
\begin{align*}
\Delta\vec{x} = \Delta\vec{x}_1 + \Delta\vec{x}_2 = (0.2,\ 0) + (173.2,\ 100) = (173.4,\ 100)
\end{align*}
with unit of meters. The second way by ``final minus initial", which is
\begin{align*}
\Delta\vec{x} = (173.2,\ 100) -(-0.2, 0) = (173.4,\ 100)
\end{align*}
with unit of meters, which is the same as above. \\

\textbf{Solution to (f).} The displacement from $t=0.8$ to $16.4$ seconds is calculated in part (e). Average velocity is a vector, calculated by that displacement over the time it takes to travel, so
\begin{align*}
\bar{\vec{v}} = (\bar{v}_x,\ \bar{v}_y) = \left( \frac{173.4}{16.4-0.54}, \ \frac{100}{16.4-0.54}\right) = \left(10.9,\ 6.3\right)
\end{align*}
with unit of meters per second.\\



\textbf{Solution to (g).} The displacement from $t = 0$ to $16.4$ seconds is, again, can be calculated in two ways. The first way by adding all the displacements, the second by ``final minus initial". Here we take the second way,
\begin{align*}
\Delta\vec{x} = (173.2,\ 100)- (0,\ 0) = (173.2,\ 100)\,,
\end{align*}
with unit of meters. Using the first way will give
\begin{align*}
\Delta\vec{x} = (-0.285,\ 0) + (0.085,\ 0) + (0.2,\ 0) + (173.2,\ 100) = (173.2,\ 100)
\end{align*}
with unit of meters, same as above. The magnitude of this vector is
\begin{align*}
|\Delta \vec{x}| = \sqrt{173.2^2 + 100^2}\,\text{m} = 200\,\text{m}\,,
\end{align*}
and the angle with respect to the horizontal is 
\begin{align*}
\theta =\arctan\left(\frac{\Delta\vec{x}_y}{\Delta \vec{x}_x}\right) =\arctan\left(\frac{100}{173.2}\right) = 30^\circ\,.
\end{align*}

\textbf{Solution to (h). } The total distance traveled is calculated by adding the distance in each segment regardless of its direction: 
\begin{enumerate}
\item First from $t=0$ to $0.2$ seconds, it moves to the left for $0.285$ meters. 
\item From $t=0.2$ to $0.54$ seconds, it moves to the right for $0.285-0.2 = 0.085$ meters.
\item  From $t=0.54$ to $0.80$ seconds, it moves to the right for $0.2-0 = 0.2$ meters.
\item From $t=0.8$ to $16.4$ seconds, it moves up-hill for $200$ meters.
\end{enumerate}
Thus in total, 
\begin{align*}
|\Delta \vec{x}| =( 0.285+0.085+0.2+200)\,\text{m} = 200.57\,\text{m}\,.
\end{align*}

\textbf{Solution to (i).} The average speed is the total distance over the time it takes to travel. With the total distance traveled calculated in part (h), 
\begin{align*}
\bar{s} = \frac{200.57}{16.4-0} \,\text{m/s} = 12.23\,\text{m/s}
\end{align*}



\end{document}


